% ============================================================
%  Kapitel 2: AleaVue -- Zufaellige Bildschirmpraesentation
% ============================================================
\chapter{Tab 1 -- AleaVue: Zufaellige Diashow}

% ============================================================
\section{Funktionsweise}
% ============================================================

AleaVue ist ein Vollbild-Zufalls-Diashow-Werkzeug.
Der Nutzer waehlt ein Verzeichnis, startet die Praesentation,
und TilVista zeigt alle enthaltenen Bilder in zufaelliger
Reihenfolge -- alle 4,5 Sekunden automatisch wechselnd.

\textbf{Ablauf Schritt fuer Schritt:}
\begin{enumerate}
  \item Nutzer setzt Verzeichnis in \klasse{DirBar}
        (Browse-Dialog oder Enter in der Pfad-Zeile)
  \item Klick auf \emph{Load from Directory} startet
        \klasse{ScanWorker} im Hintergrund-Thread --
        der GUI-Thread bleibt reaktiv
  \item Nach dem Scan stehen Bildpfade in
        \member{m\_imagePaths}, alle Pfade in
        \member{m\_allFiles}
  \item Alternativ: Doppelklick auf einen kaivo-Eintrag
        laedt gecachte Listen aus \datei{.dshow}/\datei{.catalogue}
        -- kein erneuter Scan noetig
  \item Klick auf \emph{Start AleaVue} oeffnet
        \klasse{SlideshowWindow} (Vollbild oder Fenster)
\end{enumerate}

\textbf{Tastaturkuerzel im Fenster:}
\begin{center}
\small
\begin{tabular}{ll}
\toprule
\textbf{Taste} & \textbf{Aktion} \\
\midrule
Pfeil rechts & Naechstes Bild (zufaellig) \\
Pfeil links  & Zurueck in der Anzeigehistorie \\
Leertaste    & Timer anhalten / fortsetzen \\
\textbf{S}   & Aktuelles Bild in shujuko (DB2) speichern \\
Esc          & Diashow schliessen, Sleep-Prevention zuruecksetzen \\
\bottomrule
\end{tabular}
\end{center}

% ============================================================
\section{File by File}
% ============================================================

\subsection{aleavuetab.h}

\begin{lstlisting}[style=header, caption={aleavuetab.h}]
#pragma once
#include <QStringList>
#include <QWidget>
#include <functional>

class QCheckBox;
class QLabel;
class QProgressBar;
class QPushButton;
class QThread;
class DirDatabasePanel;
class ShujukoPanel;
class SlideshowWindow;

class AleaVueTab : public QWidget
{
    Q_OBJECT
public:
    // std::function statt direkter MainWindow-Abhaengigkeit
    explicit AleaVueTab(std::function<QString()> getGlobalDir,
                        ShujukoPanel*            shujuko,
                        QWidget* parent = nullptr);

    void onDirectoryChanged(const QString& path);
    void loadFromDirectory();

    // Zugriff auf kaivo-Panel (fuer MainWindow-Signalverdrahtung)
    DirDatabasePanel* dbPanel() const { return m_dbPanel; }

signals:
    // Anfrage an MainWindow, DirBar auf diesen Pfad zu setzen
    void requestDirChange(const QString& path);

private slots:
    void onDbDirLoaded(const QString& path,
                       const QStringList& imageFiles,
                       const QStringList& allFiles);
    void onScanDone(bool ok,
                    QStringList imageFiles,
                    QStringList allFiles);
    void onScanThenOpen(bool ok,
                        QStringList imageFiles,
                        QStringList allFiles);
    void onStartSlideshow();

private:
    void buildUi();
    void scanDir(const QString& path);
    void openWindow(const QString& dir,
                    const QStringList& imagePaths);

    std::function<QString()> m_getGlobalDir;
    ShujukoPanel*            m_shujuko;  // shared, not owned

    QStringList m_imagePaths;
    QStringList m_allFiles;
    QString     m_lastScannedPath;

    QLabel*           m_lblStatus       = nullptr;
    QProgressBar*     m_pb              = nullptr;
    QPushButton*      m_btnStart        = nullptr;
    QCheckBox*        m_chkFullscreen   = nullptr;
    DirDatabasePanel* m_dbPanel         = nullptr;
    SlideshowWindow*  m_slideshowWindow = nullptr;
    QThread*          m_scanThread      = nullptr;
};
\end{lstlisting}

\subsubsection{C++-Analyse}

\paragraph{std::function als Dependency Inversion}
\verb|std::function<QString()>| ist ein Funktionsobjekt,
das jeden aufrufbaren Typ aufnehmen kann -- Lambda, Funktor,
freie Funktion. \klasse{AleaVueTab} braucht
\klasse{MainWindow} nicht zu kennen.
Das entspricht dem \emph{Dependency Inversion Principle}:
Abhaengigkeit von einer Abstraktion (Funktion), nicht von
einer konkreten Klasse.

\paragraph{ShujukoPanel* -- shared, not owned}
Der Kommentar ist semantisch wichtig.
\klasse{AleaVueTab} darf den Pointer \emph{nicht} deleten.
In modernem C++ wuerde man \verb|std::shared_ptr|
oder \verb|std::observer_ptr| verwenden; hier genuegt der
Rohzeiger, weil \klasse{MainWindow} als einziger
Eigentuemer klar definiert ist und laenger lebt als
\klasse{AleaVueTab}.

\paragraph{Zwei Scan-Slots}
\klasse{AleaVueTab} hat zwei verschiedene Slots fuer
Scan-Ergebnisse:
\begin{itemize}[noitemsep]
  \item \slot{onScanDone} -- Scan abgeschlossen, Dateiliste
    speichern, kein Fenster oeffnen (Hintergrund-Scan
    nach Verzeichniswechsel)
  \item \slot{onScanThenOpen} -- Scan abgeschlossen
    \emph{und} sofort Slideshow oeffnen (Nutzer hatte auf
    Start geklickt, bevor der Scan fertig war)
\end{itemize}

\subsubsection{Qt-Analyse}

\paragraph{DirDatabasePanel* dbPanel() const}
Getter mit \verb|const|-Qualifikation: Die Methode
veraendert kein Mitglied. \klasse{MainWindow} ruft ihn
ab, um Signal-Verbindungen zum DB-Panel herzustellen,
ohne dass \klasse{AleaVueTab} diese Verdrahtung selbst
kennen muss.

\paragraph{QProgressBar als Scan-Feedback}
\member{m\_pb} wird bei Scan-Start mit
\verb|setRange(0,100)|, \verb|setValue(0)| und
\verb|setVisible(true)| aktiviert. Fortschritt
kommt vom Signal \signal{progressChanged(int)} von
\klasse{ScanWorker} direkt per \verb|connect| an
\verb|m_pb->setValue| -- kein Code in \klasse{AleaVueTab}
noetig.

% ----------------------------------------------------------
\subsection{aleavuetab.cpp}
% ----------------------------------------------------------

\subsubsection{scanDir -- Worker-Thread starten}

\begin{lstlisting}[caption={scanDir: das Worker-Thread-Grundmuster in TilVista}]
void AleaVueTab::scanDir(const QString& path) {
    pbStart(m_pb);                 // Progressbar anzeigen
    m_btnStart->setEnabled(false); // Button sperren
    m_lastScannedPath = path;

    // Alten Thread entkoppeln (falls noch laufend)
    safeStop(m_scanThread, this);

    auto* worker = new ScanWorker(path);
    auto* thread = new QThread;
    m_scanThread = thread;

    // Worker in den neuen Thread verschieben
    worker->moveToThread(thread);

    // Verbindungen aufbauen
    connect(thread, &QThread::started,
            worker, &ScanWorker::run);
    connect(worker, &ScanWorker::progressChanged,
            m_pb,   &QProgressBar::setValue);
    connect(worker, &ScanWorker::resultReady,
            this,   &AleaVueTab::onScanDone);

    // Self-cleanup: Worker und Thread loeschen sich selbst
    connect(worker, &ScanWorker::resultReady,
            thread, &QThread::quit);
    connect(worker, &ScanWorker::resultReady,
            worker, &QObject::deleteLater);
    connect(thread, &QThread::finished,
            thread, &QObject::deleteLater);

    thread->start();
}
\end{lstlisting}

\begin{qtinfo}[title={Qt-Hintergrund: Worker-Object-Muster}]
Das \emph{Worker-Object-Muster} ist das Qt-empfohlene
Muster fuer Hintergrundarbeit ohne GUI-Blockade:
\begin{enumerate}[noitemsep]
  \item \klasse{QObject}-Subklasse (\klasse{ScanWorker})
    kapselt die Arbeit in \methode{run}
  \item \verb|moveToThread()| verschiebt den Worker
    in den Thread -- alle Slots laufen dann im Thread
  \item \signal{QThread::started} loest
    \slot{ScanWorker::run} aus
  \item Ergebnisse kommen per Signal im GUI-Thread an
    (queued connection, automatisch bei Thread-Grenze)
  \item \verb|deleteLater()| stellt thread-sicheres
    Aufraumen sicher -- nie \verb|delete| direkt aufrufen
\end{enumerate}
Niemals Qt-GUI-Methoden direkt aus einem Worker-Thread
aufrufen -- nur Signale senden!\\
Referenz: \url{https://doc.qt.io/qt-6/thread-basics.html}
\end{qtinfo}

\subsubsection{safeStop -- Schutz vor Stale-Signalen}

\begin{lstlisting}[caption={safeStop: verhindert Signale in ungueltige Slots}]
static void safeStop(QThread*& t, QObject* receiver) {
    if (!t) return;
    // Alle Verbindungen vom Thread zu diesem Objekt trennen
    t->disconnect(receiver);
    t = nullptr;
    // Der Thread laeuft weiter und raumt sich selbst auf,
    // sendet aber keine Signale mehr an "receiver".
}
\end{lstlisting}

Wenn ein neuer Scan gestartet wird, bevor der alte fertig ist,
koennten spaete Signale des alten Workers den neuen Zustand
korrumpieren. \verb|disconnect(receiver)| trennt alle
Signal-Verbindungen vom alten Thread zu \verb|this| --
ohne den Thread zu stoppen oder darauf zu warten.

\begin{warnung}[title={Achtung: Stale-Signal-Bug}]
Wuerde man \verb|safeStop| weglassen und einfach einen
zweiten Thread starten, koennte das so enden:
\begin{enumerate}[noitemsep]
  \item Scan 1 laeuft, braucht 3 Sekunden
  \item Nutzer waehlt neues Verzeichnis, Scan 2 startet
  \item Scan 2 endet nach 1 Sekunde: \verb|m_imagePaths|
    enthaelt Verzeichnis 2
  \item Scan 1 endet spaeter: \verb|m_imagePaths| wird
    mit Verzeichnis 1 ueberschrieben -- falscher Zustand!
\end{enumerate}
\end{warnung}

\subsubsection{openWindow -- SlideshowWindow erzeugen}

\begin{lstlisting}[caption={openWindow: Vollbild- oder Fenstermodus}]
void AleaVueTab::openWindow(const QString& directory,
                             const QStringList& imagePaths)
{
    const bool fullscreen = m_chkFullscreen->isChecked();
    m_slideshowWindow = new SlideshowWindow(
        directory,
        imagePaths,
        TV::logDir(),
        false,          // logErrors
        fullscreen,     // v0.5.42: Flag an Window weitergeben
        m_shujuko);

    fullscreen
        ? m_slideshowWindow->showFullScreen()
        : m_slideshowWindow->show();
}
\end{lstlisting}

\klasse{SlideshowWindow} erhaelt \member{m\_shujuko} als
Pointer, um beim S-Tastendruck direkt
\verb|m_shujuko->addBookmark(path)| aufrufen zu koennen.

% ----------------------------------------------------------
\subsection{slideshowwindow.h / slideshowwindow.cpp}
% ----------------------------------------------------------

\subsubsection{Klassendeklaration}

\begin{lstlisting}[style=header, caption={slideshowwindow.h -- Kern}]
class SlideshowWindow : public QMainWindow {
    Q_OBJECT
public:
    SlideshowWindow(const QString&     directory,
                    const QStringList& imagePaths,
                    const QString&     logPath,
                    bool               logErrors,
                    bool               fullscreen, // v0.5.42
                    ShujukoPanel*      shujuko,
                    QWidget*           parent = nullptr);
signals:
    void bookmarkAdded(const QString& path);

protected:
    void keyPressEvent(QKeyEvent* event) override;
    void resizeEvent(QResizeEvent* event) override;

private slots:
    void changeImage();

private:
    void displayImage(const QString& path);
    void goBack();
    void logError(const QString& path,
                  const QString& error);

    QString     m_directory;
    QStringList m_imagePaths;
    QString     m_logPath;
    bool        m_logErrors;
    bool        m_fullscreen = false;  // v0.5.42
    ShujukoPanel* m_shujuko;

    QLabel* m_label = nullptr;  // Anzeige-Widget
    QTimer* m_timer = nullptr;  // 4500ms Auto-Advance

    QStringList m_imageOrder;   // Anzeigeverlauf
    int         m_backtrack = 0;// < 0: in Verlauf zurueck
    QString     m_current;      // aktuell angezeigter Pfad

    int    m_showW = 1920, m_showH = 1080;
    double m_ratio = 16.0 / 9.0;
    static constexpr int kIntervalMs = 4500;
};
\end{lstlisting}

\subsubsection{C++-Analyse: Vollbild vs. Fenster (v0.5.42)}

\begin{lstlisting}[caption={Konstruktor: Vollbild- vs. Fensterlogik}]
SlideshowWindow::SlideshowWindow(...)
{
    TV::preventSleep();

    const QRect scr =
        QGuiApplication::primaryScreen()
        ->availableGeometry();

    if (m_fullscreen) {
        // Vollbild: Cursor verstecken (kein OS-Chrome)
        setCursor(Qt::BlankCursor);
        m_showW = scr.width();
        m_showH = scr.height();
    } else {
        // Fenster: Cursor sichtbar (Titelleiste braucht ihn)
        // Startet links oben bei halber Bildschirmgroesse
        m_showW = scr.width()  / 2;
        m_showH = scr.height() / 2;
        setGeometry(0, 0, m_showW, m_showH);
    }
    m_ratio = double(m_showW) / double(m_showH);
    // ...
    connect(this, &SlideshowWindow::bookmarkAdded,
            m_shujuko, &ShujukoPanel::addBookmark);

    m_timer->start();
    changeImage();
}
\end{lstlisting}

\verb|Qt::BlankCursor| versteckt den Mauszeiger vollstaendig --
nur im Vollbild sinnvoll, wo kein Fensterrahmen vorhanden ist.

\subsubsection{C++-Analyse: Zufallsreihenfolge und Rueckwaertsnavigation}

\begin{lstlisting}[caption={changeImage: Verlaufslogik mit m\_backtrack}]
void SlideshowWindow::changeImage() {
    if (m_imagePaths.isEmpty()) return;

    if (m_backtrack >= -1) {
        // Vorwaerts: neues zufaelliges Bild
        m_backtrack = 0;
        std::mt19937 rng(std::random_device{}());
        std::uniform_int_distribution<int>
            d(0, m_imagePaths.size()-1);
        const QString img = m_imagePaths.at(d(rng));
        m_imageOrder.append(img);  // in Verlauf speichern
        displayImage(img);
    } else {
        // Rueckwaerts: aus Verlauf wiederholen
        ++m_backtrack;
        const int idx =
            m_imageOrder.size() + m_backtrack;
        if (idx >= 0 && idx < m_imageOrder.size())
            displayImage(m_imageOrder.at(idx));
        else { m_backtrack = 0; changeImage(); }
    }
}
\end{lstlisting}

\member{m\_backtrack} ist ein negativer Offset in
\member{m\_imageOrder}: 0 = vorderster Eintrag,
$-1$ = eins zurueck, $-2$ = zwei zurueck usw.
Die Liste waechst vorwaerts, der Nutzer kann jederzeit
durch den Verlauf navigieren.

\begin{hinweis}
\textbf{std::mt19937 + std::random\_device:}
\verb|std::mt19937| (Mersenne Twister) ist ein
qualitativ hochwertiger Pseudozufallsgenerator --
deutlich besser als \verb|rand()|. Er wird mit einem
nicht-deterministischen Seed aus der Betriebssystem-Entropie
(\verb|std::random_device{}()|) initialisiert.
\verb|std::uniform_int_distribution<int>(a,b)| liefert
gleichverteilte Ganzzahlen in $[a, b]$ inklusive.\\
Referenz: \url{https://en.cppreference.com/w/cpp/numeric/random}
\end{hinweis}

\subsubsection{Qt-Analyse: QImageReader mit adaptiver Skalierung}

\begin{lstlisting}[caption={displayImage: effizientes Laden grosser Bilder}]
void SlideshowWindow::displayImage(const QString& path) {
    m_current = path;
    QImageReader reader(path);
    const QSize orig = reader.size();

    // Skalierung VOR dem Dekodieren setzen
    if (orig.isValid() &&
        (orig.width()  > m_showW ||
         orig.height() > m_showH))
    {
        const double s = qMin(
            double(m_showW) / orig.width(),
            double(m_showH) / orig.height());
        reader.setScaledSize(
            QSize(int(orig.width()  * s),
                  int(orig.height() * s)));
    }
    const QImage img = reader.read();
    if (img.isNull()) {
        logError(path, reader.errorString());
        changeImage();  // fehlerhaftes Bild ueberspringen
        return;
    }
    // Feinabstimmung per Seitenverhaeltnis
    QPixmap pm = QPixmap::fromImage(img);
    if (double(pm.width()) / double(pm.height()) > m_ratio)
        pm = pm.scaledToWidth(m_showW,
                              Qt::SmoothTransformation);
    else
        pm = pm.scaledToHeight(m_showH - 18,
                               Qt::SmoothTransformation);

    m_label->setPixmap(pm);
    m_timer->stop();
    m_timer->start(kIntervalMs); // Timer neu starten
    TV::preventSleep();
}
\end{lstlisting}

\begin{qtinfo}[title={Qt-Hintergrund: QImageReader vs. QPixmap(path)}]
\verb|QPixmap(path)| laedt das Bild vollstaendig in
seiner Originalaufloesung in den Speicher, bevor
skaliert wird. Bei einem 20-Megapixel-JPEG bedeutet
das: 20 MP $\times$ 4 Byte/Pixel = ca. 80 MB RAM nur
fuer ein Bild.\\
\verb|QImageReader::setScaledSize()| hingegen teilt dem
Dekoder mit, dass er \emph{beim Dekodieren} skalieren
soll -- JPEG-Dekoder koennen z.B. direkt auf 1/2, 1/4
oder 1/8 der Originalgroesse dekodieren. Das spart
erheblich Speicher und Zeit.\\
Referenz: \url{https://doc.qt.io/qt-6/qimagereader.html}
\end{qtinfo}

\subsubsection{Qt-Analyse: S-Taste und Bookmark-Signal}

\begin{lstlisting}[caption={keyPressEvent: S-Taste sendet Signal}]
void SlideshowWindow::keyPressEvent(QKeyEvent* event) {
    switch (event->key()) {
    case Qt::Key_Escape:
        setCursor(Qt::ArrowCursor);
        TV::restoreSleep();
        close();
        break;
    case Qt::Key_S:
        // emit traegt den Pfad, ShujukoPanel empfaengt ihn
        if (!m_current.isEmpty())
            emit bookmarkAdded(m_current);
        break;
    case Qt::Key_Space:
        m_timer->isActive()
            ? m_timer->stop()
            : m_timer->start(kIntervalMs);
        break;
    // ...
    default:
        QMainWindow::keyPressEvent(event); // Basisklasse!
    }
}
\end{lstlisting}

Im Konstruktor ist verbunden:
\begin{lstlisting}
connect(this,     &SlideshowWindow::bookmarkAdded,
        m_shujuko, &ShujukoPanel::addBookmark);
\end{lstlisting}

\klasse{SlideshowWindow} weiss nur, dass jemand auf das
Signal hoert -- die Implementierung von
\klasse{ShujukoPanel} ist egal. Loses Koppeln durch
Signale.

% ----------------------------------------------------------
\subsection{dirdatabasepanel.h / dirdatabasepanel.cpp (kaivo DB1)}
% ----------------------------------------------------------

\subsubsection{Aufgabe und JSON-Schema}

\klasse{DirDatabasePanel} speichert Verzeichniseintraege
mit gecachten Dateilisten. Die Datenbankdatei:

\begin{lstlisting}[language={},numbers=none,
  caption={tilvista\_dirs.json -- Struktur}]
{
  "entries": [
    {
      "name":        "Urlaub2025",
      "path":        "../../Bilder/Urlaub2025",
      "image_files": "Urlaub2025_img.dshow",
      "all_files":   "Urlaub2025_all.catalogue",
      "scanned_at":  "2025-08-01T14:23:00",
      "hidden":      false
    }
  ]
}
\end{lstlisting}

\datei{\_img.dshow} und \datei{\_all.catalogue} sind
zeilenweise Dateilisten im kaivo-Verzeichnis.
Beim Laden liest \klasse{CatalogueLoadWorker} diese
ein und liefert sie als \klasse{QStringList}.

\subsubsection{v0.5.42 Threading-Architektur}

\begin{lstlisting}[style=header,
  caption={Drei Thread-Pointer in DirDatabasePanel}]
QThread* m_thread    = nullptr; // save / index load
QThread* m_catThread = nullptr; // catalogue lesen
QThread* m_updThread = nullptr; // Verzeichnis neu scannen
\end{lstlisting}

Jede Operation laeuft in einem eigenen Thread.
Vor jeder neuen Operation wird der zugehoerige Thread
per \methode{safeStop} entkoppelt.

\begin{neuin}[title={Neu in v0.5.42: setBusy und flushPendingWrites}]
\textbf{setBusy(bool):} Waehrend ein Katalog-Load-Thread
laeuft, werden Save/Load/Delete/Update-Buttons und die
Eintragsliste deaktiviert. Das verhindert konkurrierende
Zugriffe auf \member{m\_db} (das \klasse{QJsonObject}).

\textbf{flushPendingWrites():} Beim App-Beenden pumpt
diese Methode die Event-Loop, bis der laufende Save-Thread
fertig ist:
\begin{lstlisting}
void DirDatabasePanel::flushPendingWrites() {
    if (!m_thread) return;
    QElapsedTimer t; t.start();
    while (m_thread && t.elapsed() < 3000)
        QCoreApplication::processEvents(
            QEventLoop::AllEvents, 50);
}
\end{lstlisting}
\verb|QThread::wait()| wuerde hier deadlocken:
Der Worker emittiert \signal{resultReady} als Queued
Connection zurueck in den Haupt-Thread. Dieser muss seine
Event-Loop \emph{laufen} lassen, damit das Signal
zugestellt und \verb|quit()| aufgerufen wird.
\verb|processEvents()| haelt genau genug der Event-Loop
am Leben, ohne den Thread zu blockieren.
\end{neuin}

\subsubsection{Signale von DirDatabasePanel}

\begin{lstlisting}[style=header]
signals:
    // Dateilisten geladen (aus Cache oder nach Scan)
    void dirLoaded(const QString& absPath,
                   const QStringList& imageFiles,
                   const QStringList& allFiles);
    // Aktiver Eintrag geaendert (fuer shujuko-Panel)
    void activeEntryChanged(const QString& entryName,
                             const QString& sourceDir);
    // shujuko soll Datei-Existenz pruefen
    void requestShujukoValidation();
    // Secret Mode hat sich geaendert
    void secretModeChanged(bool active);
\end{lstlisting}

\subsubsection{C++-Analyse: launchWorker -- DRY-Template-Helfer}

\begin{lstlisting}[caption={launchWorker: generische Thread-Startfunktion}]
template<typename Worker, typename Slot>
static QThread* launchWorker(Worker* worker,
                              QObject* receiver,
                              Slot slot)
{
    auto* thread = new QThread;
    worker->moveToThread(thread);
    QObject::connect(thread, &QThread::started,
                     worker, &Worker::run);
    QObject::connect(worker, &Worker::resultReady,
                     receiver, slot);
    QObject::connect(worker, &Worker::resultReady,
                     thread, &QThread::quit);
    QObject::connect(worker, &Worker::resultReady,
                     worker, &QObject::deleteLater);
    QObject::connect(thread, &QThread::finished,
                     thread, &QObject::deleteLater);
    thread->start();
    return thread;
}
\end{lstlisting}

\textbf{Warum ein Template?} Alle Worker-Klassen haben
denselben Thread-Start-Boilerplate. Das Template
(\verb|typename Worker, typename Slot|) macht die Funktion
fuer jeden Worker-Typ nutzbar ohne Code-Wiederholung
(DRY: Don't Repeat Yourself). Der Compiler instanziiert
fuer jeden Aufruf eine konkrete Version.

\subsubsection{C++-Analyse: rawName und regex}

\begin{lstlisting}[caption={rawName: Dekorations-Praefix entfernen}]
static QString rawName(const QString& displayName) {
    // Entfernt das "CIRCLE ""-Praefix von hidden entries
    static const QRegularExpression prefix("^[^ ]+ ");
    return QString(displayName).remove(prefix);
}
\end{lstlisting}

In der Eintragsliste werden versteckte Eintraege mit
dem Praefix \verb|"(kreis) "| versehen.
\verb|QRegularExpression| ist Qt's moderne Regex-Klasse
(PCRE2-basiert, Unicode-faehig). Die \verb|static|-Variable
stellt sicher, dass die Regex nur einmal kompiliert wird.

\subsubsection{Qt-Analyse: PendingCache-Struct}

\begin{lstlisting}[style=header,
  caption={PendingCache: aggregierter Zustand}]
struct PendingCache {
    QString path;
    QStringList imageFiles, allFiles;
    bool valid = false;
} m_pending;
\end{lstlisting}

Ein anonymes Struct direkt als Member-Variable --
das fasst zusammengehoerige Daten zusammen, ohne eine
benannte Klasse anlegen zu muessen. Das Flag \verb|valid|
zeigt an, ob gecachte Daten vorhanden sind, die noch
nicht in die DB geschrieben wurden.
